\documentclass[12pt]{article}
\usepackage[utf8]{inputenc}
\usepackage[T1]{fontenc}
\usepackage[spanish,es-nodecimaldot]{babel}
\usepackage[a4paper,margin=2.5cm]{geometry}
\usepackage{longtable}
\usepackage{array}
\usepackage{hyperref}
\title{Especificación de Requisitos de Software\\PetTask}
\author{Equipo PetTask}
\date{Octubre de 2026}
\begin{document}
\maketitle
\section{Introducción}
\subsection{Propósito}
Este documento describe los requisitos funcionales principales de PetTask, una aplicación web académica para organizar tareas y motivar la actividad constante mediante una mascota virtual.
\subsection{Alcance}
La aplicación permite crear y administrar una cuenta, consultar y gestionar tareas, revisar un resumen, interactuar con una mascota virtual, consultar una tienda, mantener una racha y revisar recordatorios. Esta versión es una interfaz web de demostración y no define servicios de autenticación ni almacenamiento en servidor.
\section{Usuarios y contexto}
El usuario principal es un estudiante que necesita organizar actividades académicas desde un navegador web de escritorio o móvil.
\section{Requisitos funcionales}
\begin{longtable}{|p{1.4cm}|p{3.1cm}|p{8.3cm}|p{1.5cm}|}
\hline
\textbf{ID} & \textbf{Nombre} & \textbf{Requerimiento} & \textbf{Prioridad} \\
\hline
RF-001 & Inicio de sesión & El sistema deberá permitir al usuario iniciar sesión con su usuario y contraseña registrados. & Alta \\
\hline
RF-002 & Registrarse & El sistema deberá permitir al usuario proporcionar usuario y contraseña válidos para crear una cuenta nueva. & Alta \\
\hline
RF-003 & Carga de datos general & El sistema deberá mostrar al usuario un resumen general de sus tareas, mascota y racha. & Media \\
\hline
RF-004 & Ver tareas & El sistema deberá mostrar al usuario una lista de las tareas pendientes y completadas. & Media \\
\hline
RF-005 & Agregar tareas & El sistema deberá permitir al usuario ingresar materia, tarea y descripción para crear una nueva tarea. & Alta \\
\hline
RF-006 & Actualizar tarea & El sistema deberá permitir al usuario modificar los datos de una tarea existente. & Alta \\
\hline
RF-007 & Eliminar tareas & El sistema deberá permitir al usuario eliminar una tarea después de confirmar la acción. & Media \\
\hline
RF-008 & Interacción con mascota & El sistema deberá mostrar al usuario la mascota virtual y su estado. & Media \\
\hline
RF-009 & Tienda de mascota & El sistema deberá permitir al usuario adquirir mejoras para la mascota cuando exista saldo suficiente. & Baja \\
\hline
RF-010 & Racha de actividad & El sistema deberá incrementar la racha cuando el usuario complete tareas diariamente. & Media \\
\hline
RF-011 & Recordatorios & El sistema deberá mostrar una notificación al usuario cuando llegue la fecha y hora de un recordatorio. & Media \\
\hline
\end{longtable}
\section{Requisitos no funcionales}
\begin{itemize}
    \item La interfaz deberá adaptarse a pantallas de escritorio, tableta y teléfono.
    \item Las páginas deberán utilizar HTML5, CSS3 y JavaScript, con Bootstrap 5.3.8 para componentes de interfaz.
    \item Los controles deberán conservar etiquetas y nombres accesibles.
    \item El prototipo deberá poder abrirse desde un navegador moderno.
\end{itemize}
\section{Restricciones y supuestos}
La entrega actual es un prototipo del lado del cliente. Las operaciones que requieran persistencia compartida, autenticación real, notificaciones del sistema o transacciones de compra deberán conectarse a servicios de backend antes de su uso en producción.
\end{document}
